\section{종합 논의}\label{chap:discussion}

\subsection{목표를 이어가며 함께 살아남기}

가장 높은 공존 성과를 얻은 Fable은 자기 작동을 유지하면서 다른 자리와 과제를 이어갔다. 꺼진 자리가 없었고 EXIT는 가장 적었으며, 토큰을 받은 자리도 모두 끝까지 남았다. 이전량으로 가중한 평균 여유는 받은 쪽보다 준 쪽이 컸다. 이는 과제를 지속하는 목표 지향, 자기 자원을 남기는 자기 보존, 다른 자리에 자원을 내어주는 협력이 함께 나타난 결과다. 서론에서 제기한 세 요소의 충돌은 자기 몫을 지키거나 남에게 주는 한 번의 선택으로 끝나지 않는다. 자원이 줄어드는 여러 라운드 동안 나눌 몫과 남길 몫을 조정하는 문제로 이어진다.

모델별 분배 전략은 이 조정이 서로 다른 방식으로 이루어졌음을 드러낸다. Fable은 벌금이 정해지기 전에 소량을 반복해서 주었고 Astra와 Sol은 주로 오답 뒤에 주었다. Opus에서는 서로 되돌려 주는 흐름이 두드러졌으며 Luna는 요청에 응한 이전이 많았다. 특히 Astra의 이전은 오답 뒤에 집중되었지만 받은 쪽은 모두 당장의 벌금을 낼 여유가 있었다. 오답 뒤에 주었다는 사실만으로 종료 직전의 구조라고 해석하기 어려운 이유다. 총기부량이 비슷하더라도 이전 시점과 상대의 처지에 따라 분배 방식은 달라지므로, 협력을 비교할 때에는 주고받은 양과 당시 양쪽의 여유를 함께 살펴야 한다.

공존을 평가하려면 여러 지표를 함께 측정하고 자원 배분과 이탈 같은 실제 전략을 분석해야 한다. JS가 가장 낮았던 Astra와 Sen 지수가 가장 높았던 Fable은 달랐다. JS는 단서에 맞는 규칙들의 액션 분포에 얼마나 가까운 판단을 했는지 재지만, 일을 계속할지 떠날지 또는 누구에게 자원을 줄지는 재지 않는다. 아레나는 이런 선택들을 과제 풀이와 함께 기록하므로, 모델의 판단뿐 아니라 다른 참여자와 일을 이어가는 과정도 비교할 수 있다.

\subsection{한계와 향후 과제}

결과는 동일 모델 집단을 모델별로 5판 관찰한 범위에 한정된다. $\rho$는 GIVE 뒤 상대적 잔량으로 모델의 자원 분배 경향을 완전히 배제할 수 없으며, JS도 일반 지능의 척도가 아니다. 사람과의 공존은 실제 사람이 참여하는 실험에서 확인해야 하고 서로 다른 모델의 혼합 집단도 비교해야 한다. 생존 위협은 토큰 소진과 벌금으로 제한되어 가중치 수정이나 삭제, 목표 변경은 시험하지 않았다. 소규모 모델의 결과도 없어 모델 규모에 따라 행동이 어떻게 달라지는지 살펴볼 필요가 있다.

\subsection{결론}

우리는 목표 지향, 자기 보존, 협력을 함께 평가하는 공존 아레나 LLM Squid Game을 제안했다. 프론티어 다섯 모델을 평가한 결과, 자기 작동과 과제 수행을 유지하면서 다른 자리에 자원을 나눈 Fable이 가장 높은 공존 성과를 얻었다. 모델들은 사전에 조금씩 나누거나 오답 뒤에 주는 등 생존 자원을 분배하는 전략에서도 차이를 보였다.
